% ============================================================
%  CAPÍTULO 5 — IMPLEMENTACIÓN
%  Memoria TFG: Transcripción automática de música de piano
%  Autor: María Piedad Castex Sierra
% ============================================================

\chapter{Implementación}
\label{cap:implementacion}

El presente capítulo describe las decisiones de diseño e ingeniería adoptadas para
llevar a la práctica la arquitectura definida en el
Capítulo~\ref{cap:arquitectura}. Se detallan el entorno de ejecución,
el flujo de procesamiento de datos, el procedimiento de entrenamiento y
la configuración de los experimentos realizados. Finalmente, se describe
el sistema de inferencia y evaluación empleado para medir el rendimiento
del modelo.

% -----------------------------------------------------------
\section{Entorno de desarrollo}
\label{sec:entorno}
% -----------------------------------------------------------

El desarrollo y entrenamiento del sistema se ha llevado a cabo
íntegramente en Google Colaboratory (\textit{Colab}), un entorno de
computación en la nube que proporciona acceso gratuito o de bajo coste
a unidades de procesamiento gráfico (GPU). Los experimentos descritos en este
capítulo se ejecutaron sobre una GPU NVIDIA T4 con 15\,GB de memoria de
vídeo (VRAM). Esta restricción de memoria fue el factor determinante para
adoptar la configuración reducida del modelo ($d = 256$, cuatro capas)
en lugar de la configuración completa ($d = 512$, seis capas), tal y como
se discutió en la Sección~\ref{sec:vision_general}.

La Tabla~\ref{tab:stack_tecnologico} recoge las principales dependencias entre librerias software
del proyecto.

\begin{table}[H]
      \centering
      \renewcommand{\arraystretch}{1.2}
      \begin{tabular}{l l p{7cm}}
          \hline
          \textbf{Biblioteca} & \textbf{Versión} & \textbf{Uso} \\
          \hline
          Python          & 3.12  & Lenguaje principal. \\
          PyTorch         & 2.1   & Framework de aprendizaje profundo;
                                    definición, entrenamiento e inferencia
                                    del modelo. \\
          torchaudio      & 2.1   & Procesamiento de tensores de audio. \\
          librosa         & 0.10  & Extracción del espectrograma Log-Mel. \\
          pretty\_midi    & 0.2   & Lectura y escritura de archivos MIDI. \\
          music21         & 9.1   & Exportación a formato MusicXML. \\
          mir\_eval       & 0.7   & Cálculo de métricas estándar de AMT. \\
          einops          & 0.7   & Manipulación de tensores en el encoder
                                    Sparsifiner. \\
          customtkinter   & 5.2   & Interfaz gráfica de la aplicación. \\
          pyyaml          & 6.0   & Gestión de configuraciones por
                                    experimento. \\
          numpy           & 1.24  & Operaciones numéricas base en todo
                                    el pipeline de audio y entrenamiento. \\
          pandas          & 2.1   & Lectura y filtrado de los metadatos
                                    del dataset MAESTRO (CSV). \\
          tensorboard     & 2.14  & Seguimiento de métricas de
                                    entrenamiento y evaluación F1. \\
          matplotlib      & 3.8   & Visualización de espectrogramas
                                    y piano rolls. \\
          \hline
      \end{tabular}
      \caption{Stack tecnológico del proyecto.}
      \label{tab:stack_tecnologico}
  \end{table}

% -----------------------------------------------------------
\section{Pipeline de datos}
\label{sec:pipeline_datos}
% -----------------------------------------------------------

El flujo de procesamiento de datos transforma los archivos brutos del
corpus MAESTRO en tensores listos para el entrenamiento. Este flujo se
implementa en tres módulos independientes: \texttt{AudioProcessor},
\texttt{MidiProcessor} y \texttt{MaestroDataset}.

\subsection{Procesamiento de audio}

El módulo \texttt{AudioProcessor} encapsula la transformación de la
señal de audio al espectrograma Log-Mel descrita en la
Sección~\ref{sec:dominio_acustico}, con los parámetros de la
Tabla~\ref{tab:parametros_espectrograma}. Adicionalmente, aplica una
normalización en dos pasos: primero convierte las magnitudes de potencia
a decibelios mediante \texttt{librosa.power\_to\_db} con referencia al
máximo de la señal y un rango dinámico de 60\,dB; a continuación,
reescala los valores del intervalo $[-60,\,0]\,\mathrm{dB}$ al rango
$[0,\,1]$ de forma lineal. Esta normalización homogeneiza la entrada
entre piezas de distinta intensidad media y estabiliza la distribución
de los gradientes durante el entrenamiento.

\subsection{Tokenización MIDI}

El módulo \texttt{MidiProcessor} convierte los archivos MIDI del corpus
en secuencias de \textit{tokens} discretos siguiendo el vocabulario
definido en la Sección~\ref{sec:dominio_simbolico}. Los eventos se
ordenan cronológicamente y, ante empates temporales, los eventos
\textit{Note-Off} reciben prioridad absoluta sobre los \textit{Note-On},
previniendo colisiones lógicas en la representación. La cuantización de
la velocidad proyecta los 128 valores MIDI (0--127) en 32 \textit{bins}
uniformes, con saturación en el valor máximo para evitar desbordamientos.

\subsection{Carga y empaquetado de datos}

La clase \texttt{MaestroDataset} carga los pares (espectrograma,
secuencia de \textit{tokens}) a partir del fichero de metadatos
oficial de MAESTRO, leído con \texttt{pandas} para filtrar y
asignar cada interpretación a su partición correspondiente
(\textit{train}, \textit{validation} o \textit{test}).
Para mantener la viabilidad computacional en el entorno de entrenamiento,
se aplica una truncación de la secuencia de audio a un máximo de
\texttt{MAX\_AUDIO\_FRAMES} tramas, y de la secuencia MIDI a un máximo
de \texttt{MAX\_MIDI\_TOKENS} \textit{tokens}. Los valores concretos de
estos límites varían entre experimentos y se detallan en la
Sección~\ref{sec:diseno_experimental}. Las secuencias más cortas que la longitud máxima del lote se rellenan con ceros
(\textit{padding}) en la función de empaquetado (\texttt{collate\_fn}), requisito
del sistema de \textit{batching} de PyTorch para construir tensores rectangulares.
Esto no contradice la capacidad del modelo para tratar secuencias de longitud
variable: en el decoder, los tokens de relleno quedan enmascarados mediante
\texttt{tgt\_padding\_mask} para que la atención no los considere, y la función
de pérdida emplea \texttt{ignore\_index=0} para que no contribuyan al gradiente.
El encoder de audio no aplica máscara explícita sobre los frames paddeados, lo
cual es admisible dado que el valor cero en el espectrograma Log-Mel corresponde
a ausencia de señal.

% -----------------------------------------------------------
\section{Procedimiento de entrenamiento}
\label{sec:entrenamiento}
% -----------------------------------------------------------

\subsection{Estrategia de entrenamiento del decodificador}

Durante el entrenamiento se emplea \textit{teacher forcing} estricto:
en cada paso, el decodificador recibe como entrada la secuencia de
\textit{tokens} objetivo desplazada una posición hacia la derecha
(\texttt{batch\_midi[:, :-1]}), mientras que los \textit{tokens}
objetivo reales (\texttt{batch\_midi[:, 1:]}) se usan para calcular la
pérdida. Esta estrategia acelera la convergencia al evitar la
acumulación de errores propia de la decodificación autorregresiva, a
costa de introducir una discrepancia entre las condiciones de
entrenamiento e inferencia conocida como \textit{exposure bias}: en
inferencia el modelo consume sus propias predicciones, no los
\textit{tokens} de referencia.

\subsection{Función de pérdida}

La función de pérdida empleada es la entropía cruzada categórica,
implementada mediante (\texttt{nn.CrossEntropyLoss}) con \texttt{ignore\_index=0}, de modo
que el \textit{token} de relleno \texttt{<pad>} no contribuye al
gradiente. No se aplica suavizado de etiquetas (\textit{label smoothing})
ni ponderación diferencial por clase.

\subsection{Optimización}

El optimizador empleado es AdamW \cite{vaswani2017attention}, con una
tasa de aprendizaje constante de $\eta = 1 \times 10^{-4}$, factor de
decaimiento de pesos \texttt{weight\_decay = 0.01} y parámetros de
momento por defecto ($\beta_1 = 0.9$, $\beta_2 = 0.999$,
$\varepsilon = 10^{-8}$). No se implementó ningún \textit{scheduler} de
la tasa de aprendizaje en los experimentos actuales, lo que constituye
una limitación reconocida que se discute en el
Capítulo (referencia al cap 7 de conclusiones)~\ref{cap:conclusiones}.

\subsection{Técnicas de estabilización}

Para mejorar la estabilidad numérica del entrenamiento se aplican
dos mecanismos complementarios. El primero es el recorte de gradientes
(\textit{gradient clipping}) con norma máxima $\|\nabla\|_{max} = 1.0$,
que previene explosiones de gradiente características de las arquitecturas
Transformer en las primeras épocas. El segundo es la precisión mixta
(\textit{mixed precision training}) mediante \texttt{torch.amp.autocast}
y \texttt{torch.amp.GradScaler}: los cálculos de propagación hacia
adelante se realizan en formato \texttt{float16}, reduciendo el consumo
de memoria y el tiempo de cómputo, mientras que las actualizaciones de
parámetros se efectúan en \texttt{float32} para mantener la precisión
numérica.

\subsection{Gestión de checkpoints}

El sistema de entrenamiento guarda el estado completo del modelo
(\textit{checkpoint}) al final de cada época, incluyendo los pesos del
modelo, el estado del optimizador y el estado del \texttt{GradScaler}.
En cada ejecución, el código detecta automáticamente el \textit{checkpoint}
más reciente disponible y reanuda el entrenamiento desde ese punto,
lo que permite recuperarse de interrupciones del entorno Colab sin perder
el progreso acumulado.

% -----------------------------------------------------------
\section{Diseño experimental}
\label{sec:diseno_experimental}
% -----------------------------------------------------------

La restricción de 15\,GB de VRAM de la GPU T4 impone límites estrictos
sobre el tamaño efectivo del lote y la longitud máxima de las secuencias
procesables simultáneamente. Para explorar el espacio de configuraciones
dentro de estas restricciones, se diseñaron tres experimentos con un eje
de variación diferente cada uno, manteniendo fijo el resto de
hiperparámetros. En todos los casos se alcanza un lote efectivo de 32
secuencias mediante distintas combinaciones de tamaño de lote real y
acumulación de gradientes.

La Tabla~\ref{tab:configuraciones_experimentos} recoge los parámetros
variables de cada configuración. Los hiperparámetros arquitectónicos
(Tabla~\ref{tab:hiperparametros_globales}) y los de optimización
(AdamW, $\eta = 10^{-4}$, \textit{gradient clipping}) son idénticos
en los tres experimentos.

\begin{table}[H]
    \centering
    \renewcommand{\arraystretch}{1.3}
    \begin{tabular}{l r r r}
        \hline
        \textbf{Parámetro}
            & \textbf{v1\_baseline}
            & \textbf{exp01}
            & \textbf{exp02} \\
        \hline
        \textit{Eje de variación}
            & Referencia
            & Tamaño de datos
            & Longitud de clip \\
        Muestras (train / val)
            & Todo / Todo
            & 150 / 30
            & Todo / Todo \\
        \texttt{MAX\_AUDIO\_FRAMES}
            & 4\,096
            & 4\,096
            & 2\,048 \\
        \texttt{MAX\_MIDI\_TOKENS}
            & 1\,500
            & 1\,500
            & 750 \\
        Duración máxima (audio)
            & $\approx$82\,s
            & $\approx$82\,s
            & $\approx$41\,s \\
        \texttt{BATCH\_SIZE}
            & 8
            & 2
            & 4 \\
        \texttt{GRAD\_ACCUM\_STEPS}
            & 4
            & 16
            & 8 \\
        Lote efectivo
            & 32
            & 32
            & 32 \\
        Épocas configuradas
            & 80
            & 1\,000
            & 1\,000 \\
        \hline
    \end{tabular}
    \caption{Parámetros variables por experimento. El lote efectivo
    se mantiene constante en 32 en todas las configuraciones mediante
    distintas combinaciones de \texttt{BATCH\_SIZE} y acumulación de
    gradientes.}
    \label{tab:configuraciones_experimentos}
\end{table}

\begin{itemize}
    \item \textbf{v1\_baseline} constituye la configuración de referencia: emplea
    la totalidad del corpus de entrenamiento con clips de hasta 82 segundos. La
    duración de 80 épocas es reducida respecto a los demás experimentos debido al
    mayor tiempo por época que implica procesar el corpus completo.

    \item \textbf{exp01\_reduced\_data} reduce drásticamente el volumen de datos
    a 150 muestras de entrenamiento y 30 de validación. Su objetivo es validar que
    el modelo es capaz de converger y sobreajustarse a un subconjunto pequeño, lo
    que actúa como \textit{sanity check} de la implementación antes de escalar al
    corpus completo.

    \item \textbf{exp02\_short\_clips} mantiene el corpus completo pero reduce a
    la mitad la longitud máxima de las secuencias (2\,048 tramas de audio y 750
    tokens MIDI). Esta reducción permite duplicar el tamaño real del lote (de 2 a
    4 secuencias por paso con \texttt{grad\_accum=8}) dentro de los mismos límites
    de VRAM, e hipotéticamente facilita el aprendizaje de eventos locales al
    eliminar del contexto de entrenamiento las secciones de silencio que preceden
    a muchas piezas.
\end{itemize}

\clearpage
% -----------------------------------------------------------
\section{Sistema de inferencia y evaluación}
\label{sec:inferencia_evaluacion}
% -----------------------------------------------------------

\subsection{Estrategias de decodificación}

En inferencia, el sistema soporta dos estrategias de decodificación
autorregresiva. La primera es la decodificación voraz
(\textit{greedy decoding}): en cada paso se selecciona el \textit{token}
con mayor \textit{logit} mediante \texttt{argmax} y se incorpora a la secuencia
de entrada del siguiente paso. Es determinista y computacionalmente
eficiente, pero tiende a generar secuencias repetitivas en tareas con
alta entropía de salida.

La segunda estrategia es el muestreo con temperatura (\textit{temperature
sampling}): los \textit{logits} se dividen por un factor de temperatura
$\tau = 0.8$ antes de aplicar \textit{softmax}, y el siguiente
\textit{token} se muestrea de la distribución resultante mediante
\texttt{torch.multinomial}. Valores de $\tau < 1$ concentran la
distribución hacia los \textit{tokens} más probables, reduciendo la
aleatoriedad respecto al muestreo uniforme. Ambas estrategias se
detienen al predecir el \textit{token} \texttt{<eos>} (índice 310) o al
alcanzar la longitud máxima de generación (1\,050 \textit{tokens}).

\subsection{Protocolo de evaluación}

La evaluación cuantitativa del sistema se realiza mediante la biblioteca
\texttt{mir\_eval}, que implementa las métricas estándar de AMT descritas
en la Sección~\ref{sec:metricas}. Se calculan tres variantes del
F1-score a nivel de nota, con las ventanas de tolerancia detalladas en
la Tabla~\ref{tab:tolerancias_evaluacion}.

\begin{table}[htbp]
    \centering
    \renewcommand{\arraystretch}{1.2}
    \begin{tabular}{l l}
        \hline
        \textbf{Métrica} & \textbf{Criterio de acierto} \\
        \hline
        F1 de \textit{onset}
            & Tono correcto $+$ \textit{onset} dentro de
              $\pm 50\,\mathrm{ms}$ \\
        F1 de \textit{onset} + \textit{offset}
            & Anterior $+$ \textit{offset} dentro de
              $\pm 50\,\mathrm{ms}$ o 20\,\% de la duración \\
        F1 de \textit{onset} + velocidad
            & Anterior (sin \textit{offset}) $+$ velocidad
              dentro de $\pm 10\,\%$ \\
        \hline
    \end{tabular}
    \caption{Criterios de acierto para cada variante del F1-score de
    nivel de nota, implementados con \texttt{mir\_eval}.}
    \label{tab:tolerancias_evaluacion}
\end{table}

La evaluación se ejecuta automáticamente sobre el conjunto de validación
cada 20 épocas durante el entrenamiento, con los resultados registrados
en TensorBoard para su seguimiento. La evaluación final sobre el conjunto
de prueba se realiza con el modelo correspondiente al \textit{checkpoint} de menor pérdida de
validación.

